\itemhead{第1章 小结与复习 · 第2题（教材p.28）}{XJ-B1-C01-SU-P028-Y1103}
\label{XJ-B1-C01-SU-P028-Y1103}
{\footnotesize 湘教版教材p.28 / PDF第35页 / 第1章 小结与复习 / 第2题}
\field{原题}{集合的基本概念，可以浓缩为一个符号“$\in$”，有关集合的其他概念和符号都可以用符号“$\in$”来刻画。试结合集合的包含关系和集合的交、并、补运算，体会如何用“$\in$”来刻画这些关系与运算。}
\field{参考答案}{$A\subseteq B\Leftrightarrow\forall x(x\in A\Rightarrow x\in B)$；$x\in A\cap B\Leftrightarrow x\in A$ 且 $x\in B$；$x\in A\cup B\Leftrightarrow x\in A$ 或 $x\in B$；$x\in\complement_U A\Leftrightarrow x\in U$ 且 $x\notin A$。}
{\small \field{关联知识}{S03 子集、真子集与集合相等（教材帮p.11、12、13 / PDF20、21、22）；S04 交集、并集及解集的联系（教材帮p.19、20 / PDF28、29）；S05 全集、补集与德摩根律（教材帮p.21 / PDF30）}
\field{方法与易错点}{用逐元素检验或数轴包含；参数题先确认是否允许空集，再讨论端点。 易错：遗漏空集；把元素个数相等当作集合相等；把子集符号当作属于符号。}
\field{答案依据}{答案册未提供本栏目的独立对应答案；依据原题、教材概念补写参考回应。开放性总结不设唯一数值答案。}}
